% GENERATED FILE. Edit canonical lessons, then run npm run generate:books.
% canonical-source-hash: fnv1a64:06194ddf40b8dfa3
% canonical-lessons: MR-C194-phon-karne, MR-C194-carca-karne, MR-C194-manai-karne, MR-C194-raksan-karne, MR-C194-marne2

\chapter{To phone, To discuss, To forbid, To protect, To die}
\label{ch:mr-to-phone-to-discuss-to-forbid-to-protect-to-die}
\hlchaptermodality{194}

\begin{chapteropening}
\textbf{By the end of this chapter:} \emph{I can say to phone, to discuss, to forbid, to protect and to die.}

Complete the last lesson of chapter 194: I can say to phone, to discuss, to forbid, to protect and to die.
\end{chapteropening}

\section[\texorpdfstring{\mr{फोन} \mr{करणे} (phon karṇe)}{फोन करणे (phon karṇe)}]{\mr{फोन} \mr{करणे} (phon karṇe) — to phone}

\label{lesson:MR-C194-phon-karne}



\begin{quote}
Before the new one: say the Marathi for to dry, then the Marathi for to iron.
\end{quote}

\subsection*{You'll want to know: \mr{फोन} \mr{करणे}}
\textbf{\mr{फोन} \mr{करणे}} — \emph{phon karṇe} — \textquotedblleft{}to phone\textquotedblright{}.

\begin{grammarlens}[title={What you've built}]
One more word for this chapter.
\end{grammarlens}

\subsection*{Your turn}
\begin{itemize}
\raggedright
  \item \emph{Say it:} \emph{phon karṇe}
  \item \emph{Say it:} \emph{phon karṇe}, once more
  \item \emph{Say it:} say \emph{phon karṇe}
  \item \emph{Recall:} say the Marathi for to dry, then the Marathi for to iron, then say \emph{phon karṇe} again
\end{itemize}

\begin{tcolorbox}[breakable,colback=teal!4,colframe=teal!35!black,title={Before you move on}]
What does \mr{फोन} \mr{करणे} mean? (\textquotedblleft{}To phone\textquotedblright{}.) Say it once more.
\end{tcolorbox}
\section[\texorpdfstring{\mr{चर्चा} \mr{करणे} (carcā karṇe)}{चर्चा करणे (carcā karṇe)}]{\mr{चर्चा} \mr{करणे} (carcā karṇe) — to discuss}

\label{lesson:MR-C194-carca-karne}



\begin{quote}
Before the new one: say the Marathi for to iron, then the Marathi for to phone.
\end{quote}

\subsection*{You'll want to know: \mr{चर्चा} \mr{करणे}}
\textbf{\mr{चर्चा} \mr{करणे}} — \emph{carcā karṇe} — \textquotedblleft{}to discuss\textquotedblright{}.

\begin{grammarlens}[title={What you've built}]
One more word for this chapter.
\end{grammarlens}

\subsection*{Your turn}
\begin{itemize}
\raggedright
  \item \emph{Say it:} \emph{carcā karṇe}
  \item \emph{Say it:} \emph{carcā karṇe}, once more
  \item \emph{Say it:} say \emph{carcā karṇe}
  \item \emph{Recall:} say the Marathi for to iron, then the Marathi for to phone, then say \emph{carcā karṇe} again
\end{itemize}

\begin{tcolorbox}[breakable,colback=teal!4,colframe=teal!35!black,title={Before you move on}]
What does \mr{चर्चा} \mr{करणे} mean? (\textquotedblleft{}To discuss\textquotedblright{}.) Say it once more.
\end{tcolorbox}
\section[\texorpdfstring{\mr{मनाई} \mr{करणे} (manāī karṇe)}{मनाई करणे (manāī karṇe)}]{\mr{मनाई} \mr{करणे} (manāī karṇe) — to forbid}

\label{lesson:MR-C194-manai-karne}



\begin{quote}
Before the new one: say the Marathi for to phone, then the Marathi for to discuss.
\end{quote}

\subsection*{You'll want to know: \mr{मनाई} \mr{करणे}}
\textbf{\mr{मनाई} \mr{करणे}} — \emph{manāī karṇe} — \textquotedblleft{}to forbid\textquotedblright{}.

\begin{grammarlens}[title={What you've built}]
One more word for this chapter.
\end{grammarlens}

\subsection*{Your turn}
\begin{itemize}
\raggedright
  \item \emph{Say it:} \emph{manāī karṇe}
  \item \emph{Say it:} \emph{manāī karṇe}, once more
  \item \emph{Say it:} say \emph{manāī karṇe}
  \item \emph{Recall:} say the Marathi for to phone, then the Marathi for to discuss, then say \emph{manāī karṇe} again
\end{itemize}

\begin{tcolorbox}[breakable,colback=teal!4,colframe=teal!35!black,title={Before you move on}]
What does \mr{मनाई} \mr{करणे} mean? (\textquotedblleft{}To forbid\textquotedblright{}.) Say it once more.
\end{tcolorbox}
\section[\texorpdfstring{\mr{रक्षण} \mr{करणे} (rakṣaṇ karṇe)}{रक्षण करणे (rakṣaṇ karṇe)}]{\mr{रक्षण} \mr{करणे} (rakṣaṇ karṇe) — to protect}

\label{lesson:MR-C194-raksan-karne}



\begin{quote}
Before the new one: say the Marathi for to discuss, then the Marathi for to forbid.
\end{quote}

\subsection*{You'll want to know: \mr{रक्षण} \mr{करणे}}
\textbf{\mr{रक्षण} \mr{करणे}} — \emph{rakṣaṇ karṇe} — \textquotedblleft{}to protect\textquotedblright{}.

\begin{grammarlens}[title={What you've built}]
One more word for this chapter.
\end{grammarlens}

\subsection*{Your turn}
\begin{itemize}
\raggedright
  \item \emph{Say it:} \emph{rakṣaṇ karṇe}
  \item \emph{Say it:} \emph{rakṣaṇ karṇe}, once more
  \item \emph{Say it:} say \emph{rakṣaṇ karṇe}
  \item \emph{Recall:} say the Marathi for to discuss, then the Marathi for to forbid, then say \emph{rakṣaṇ karṇe} again
\end{itemize}

\begin{tcolorbox}[breakable,colback=teal!4,colframe=teal!35!black,title={Before you move on}]
What does \mr{रक्षण} \mr{करणे} mean? (\textquotedblleft{}To protect\textquotedblright{}.) Say it once more.
\end{tcolorbox}
\section[\texorpdfstring{\mr{मरणे} (marṇe)}{मरणे (marṇe)}]{\mr{मरणे} (marṇe) — to die}

\label{lesson:MR-C194-marne2}



\begin{quote}
Before the new one: say the Marathi for to forbid, then the Marathi for to protect.
\end{quote}

\subsection*{You'll want to know: \mr{मरणे}}
\textbf{\mr{मरणे}} — \emph{marṇe} — \textquotedblleft{}to die\textquotedblright{}.

\begin{grammarlens}[title={What you've built}]
One more word for this chapter.
\end{grammarlens}

\subsection*{Your turn}
\begin{itemize}
\raggedright
  \item \emph{Say it:} \emph{marṇe}
  \item \emph{Say it:} \emph{marṇe}, once more
  \item \emph{Say it:} say \emph{marṇe}
  \item \emph{Recall:} say the Marathi for to forbid, then the Marathi for to protect, then say \emph{marṇe} again
\end{itemize}

\begin{tcolorbox}[breakable,colback=teal!4,colframe=teal!35!black,title={Before you move on}]
What does \mr{मरणे} mean? (\textquotedblleft{}To die\textquotedblright{}.) Say it once more.
\end{tcolorbox}
